\section{\label{I_3_A}Indexation} 

L’indexation documentaire peut être définie de la manière suivante :  

\begin{quote} 
	
	"L'indexation est cette technique consistant à caractériser le contenu d'un document et l'information qu'il détient de manière à le retrouver quand on effectue des recherches sur l'un des sujets dont il traite. [...] L'indexation permet de s'orienter dans la masse des documents et d'organiser ses connaissances. "\footcite{zotero-item-582} 
	
\end{quote} 

Afin de "fournir un point d'accès sûr et vérifié à l'information", l'indexation implique la mise en relation des termes produits par la description avec des référentiels, des thésaurus et des listes d'autorité, afin de pouvoir les intégrer au modèle de données.\footcite{zotero-item-675} Ainsi, par exemple, le réalisateur d’un film est indexé au niveau Œuvre, par une liste contrôlée d’autorités prévue par l’outil de gestion des collections.  

La structuration des métadonnées d’indexation permet d'améliorer la découvrabilité de ces métadonnées via des outils de recherche et la navigation entre les différentes connaissances produites.\footcite{zotero-item-675} Une indexation reliée à des thésaurus contrôlés permet notamment à l’utilisateur d’effectuer des requêtes croisées complexes au sein de l’outil de gestion des collections. Cette tâche, prévue par le logiciel, effectue le lien indispensable entre la description produite et l’exploitabilité des données par l’utilisateur via des outils de recherche. À terme, l'indexation permet l’interopérabilité des données et leur partage sur le web. \footcite{zotero-item-675}  

Toutefois, la nature des archives audiovisuelles complexifie l’indexation pour l’utilisateur. En effet, nous l’avons évoqué, de telles archives se constituent comme des “boites noires”, nécessitant le matériel de lecture adéquat afin d’en extraire le sens et de pouvoir l’indexer. De plus, l’indexation du corpus que nous avons décrit, composé de rushes documentaires, suppose en amont la division du flux vidéo en différentes unités de sens thématiques. La description ne peut pas s’appliquer de manière précise dans le cas d’un document vidéo brut contenant de multiples séquences spatio-temporelles. Dans le cas de ce type d’archives audiovisuelles, la segmentation temporelle constituerait alors la condition préalable à toute opération d’indexation.